\section{Estimate of the eight length} \label{app8}
Let $D:=2\Delta\theta/\Delta\varphi$, $s:=\sin\theta_0$, $\alpha:=2\Delta\theta=D\,\Delta\varphi$, $\beta:=\Delta\varphi\,s$; recall that % foo
$$L_{81}=\int_0^{2\pi}\sqrt{[2\Delta\theta\cos2\tau]^2+\sin^2(\theta_0+\Delta\theta\sin2\tau)\,[\Delta\varphi\cos\tau]^2}\,\dd\tau,$$
and let $\widetilde L_{81}$ be the quantity obtained by replacing $\sin^2(\theta_0+\Delta\theta\sin2\tau)$ by $\sin^2\theta_0$:
$$\widetilde L_{81}:=\int_0^{2\pi}\sqrt{\alpha^2\cos^22\tau+\beta^2\cos^2\tau}\,\dd\tau .$$
The error of this approximation is analysed as follows, with $x:=\Delta\theta$ and $E(x):=L_{81}-\widetilde L_{81}$. Decoupling the coefficient of $\cos2\tau$ from the amplitude in the sine, put, for $a\in\mathbf{R}$ and $\varepsilon$ small (only $a^2$ enters, so $\Lambda$ is even in $a$),
$$\Lambda(a,\varepsilon):=\int_0^{2\pi}\sqrt{a^2\cos^22\tau+\Delta\varphi^2\cos^2\tau\,\sin^2(\theta_0+\varepsilon\sin2\tau)}\,\dd\tau,$$
so that $L_{81}=\Lambda(2x,x)$, $\widetilde L_{81}=\Lambda(2x,0)$ and $E(x)=\Lambda(2x,x)-\Lambda(2x,0)$. Then $E$ is twice differentiable at $0$, with $E(0)=0$ and $E'(0)=0$, the latter because $\partial_\varepsilon\Lambda(a,0)=0$ for every $a\in\mathbf{R}$ (the integrand of $\partial_\varepsilon\Lambda(a,0)$ is odd under $\tau\mapsto\pi-\tau$ whereas the radicand is even), and $E''(0)=\Lambda_{\varepsilon\varepsilon}(0,0)$. The latter is computable in closed form because at $a=0$ the integrand collapses to $\Delta\varphi\,|\cos\tau|\sin(\theta_0+\varepsilon\sin2\tau)$, whose second derivative at $\varepsilon=0$ is $-\Delta\varphi\,s\,|\cos\tau|\sin^22\tau$, whence
$$E''(0)=-\Delta\varphi\,s\int_0^{2\pi}|\cos\tau|\sin^22\tau\,\dd\tau=-\frac{32}{15}\,\Delta\varphi\,s\ne0 .$$
The Taylor--Young formula $E(x)=\tfrac12E''(0)x^2+o(x^2)$ therefore gives
$$L_{81}=\widetilde L_{81}-\frac{16}{15}\,\Delta\varphi\,s\,\Delta\theta^2+o(\Delta\theta^2),$$
and the approximation is of order exactly two. Set
$$C:=\frac{\alpha^2}{\alpha^2+\beta^2}=\frac{D^2}{D^2+s^2}\in(0,1).$$
By the symmetries of the integrand, and with $S:=\sin^2\tau$ (so $\dd S=2\sqrt{S(1-S)}\,\dd\tau$ and $\cos2\tau=1-2S$),
$$\begin{gathered}\widetilde L_{81}=4\alpha\int_0^1\frac{N(S)\,\dd S}{w},\qquad w^2:=S(1-S)\,N(S),\\ N(S):=S^2-\sigma S+\varpi,\qquad \sigma:=\frac{1+3C}{4C},\qquad \varpi:=\frac1{4C},\end{gathered}$$
where $N>0$ on $[0,1]$, since $N(0)=\varpi$ and $N(1)=1/4$. The curve $w^2=S(1-S)N(S)$ is a quartic with four simple roots for $C\ne1/9$, hence of genus one, and the limits $0$ and $1$ are two of its roots, so $\widetilde L_{81}$ is a complete elliptic integral. Sending the root $S=0$ to infinity with $t:=1/S$, and removing the quadratic term with
$$x:=\frac{t-\tfrac13(2+3C)}{16\,C},\qquad y:=\frac{w}{16\,C\,S^2},$$
the curve becomes $y^2=4x^3-g_2x-g_3$ with
$$g_2:=\frac{9C^2-9C+1}{192\,C^2},\qquad g_3:=\frac{(3C-1)(18C^2-21C+2)}{27648\,C^3},\qquad g_2^3-27g_3^2=\frac{(C-1)(9C-1)}{1048576\,C^4},$$
which are rational in $C$, hence in $D^2$ and $s^2$. The root $S=1$ becomes $x=e_r:=(1-3C)/(48C)$, and the point $S=\infty$ becomes $x=x_0:=-(2+3C)/(48C)$. Let $\wp$ be the Weierstrass function of $(g_2,g_3)$, with period lattice generated by $2\omega$ and $2\omega'$, and let $\zeta$ be defined by $\zeta'=-\wp$ and $\zeta(z)-1/z\to0$ as $z\to0$, so that $\zeta(z+2\omega)=\zeta(z)+2\eta$. Parametrising $x=\wp(u)$, $y=-\wp'(u)$, one has $\dd S/w=\dd u$, and $S\in[0,1]$ corresponds to $u\in[0,\omega]$, where
$$\begin{gathered}\omega:=\int_{e_r}^{\infty}\frac{\dd x}{\sqrt{4x^3-g_2x-g_3}},\qquad \eta:=\zeta(\omega),\\ b:=\int_{-\infty}^{x_0}\frac{\dd x}{\sqrt{g_3+g_2x-4x^3}},\qquad a:=-ib,\end{gathered}$$
so that $\omega$ is the real \emph{half}-period ($\wp(\omega)=e_r$ and $\wp'(\omega)=0$), and $\wp(a)=x_0$, $\wp'(a)=i/(16C)$. The lattice is rectangular iff $g_2^3-27g_3^2>0$, i.e.\ iff $C<1/9$ (equivalently $s>2\sqrt2\,D$), in which case $\omega'$ is purely imaginary. It is rhombic (non-rectangular) for $1/9<C<1$, and degenerate for $C=1/9$. In terms of $u$, with $t=16C(\wp(u)-x_0)$,
$$\frac{N\,\dd S}{w}=\Bigl(\varpi-\frac{\sigma}{t}+\frac1{t^2}\Bigr)\dd u,$$
and the classical identity $\wp'(a)/(\wp(u)-\wp(a))=\zeta(u-a)-\zeta(u+a)+2\zeta(a)$ gives
$$\int_0^\omega\frac{\dd u}{\wp(u)-\wp(a)}=\frac{2\bigl(\omega\,\zeta(a)-\eta\,a\bigr)-i\pi}{\wp'(a)},$$
where the constant $-i\pi$ is the branch fixed by the behaviour for $b\to0^+$, and the integral with a squared denominator is obtained by differentiating with respect to $a$, using $\wp'(a)^2=-1/(256C^2)$ and $\wp''(a)=(1+7C)/(128C^2)$. Collecting the three kinds of terms, only $\omega$, $\eta$ and $\zeta(a)$ remain:
$$\widetilde L_{81}=4\alpha\,I(C),\qquad I(C):=\frac{(4-3C)\,\omega}{24\,C}+2\eta+\frac{1-C}{8\,C}\Bigl[\pi+2i\bigl(\omega\,\zeta(a)-\eta\,a\bigr)\Bigr],$$
which is real since $a\in i\mathbf{R}$. The third-kind term has coefficient $(1-C)/(8C)$, which vanishes only at $C=1$, and the limit $C\to1$ ($\theta_0\to0$ or $\Delta\varphi\to0$) gives $I\to1$, i.e.\ $\widetilde L_{81}\to4\alpha=8\Delta\theta$. For small $D$ (i.e.\ $C$, that is $\Delta\theta$, small) the four terms of $I(C)$ are individually of order $1/C$ and cancel down to
$$I(C)=\frac1{\sqrt C}\Bigl[1+\frac C4\Bigl(\ln\frac{16}{C}-\frac{11}{3}\Bigr)+O(C^2\ln C)\Bigr],$$
equivalently $\widetilde L_{81}=4\rho\bigl[1+\tfrac C4(\ln\tfrac{16}{C}-\tfrac{11}3)+\dots\bigr]$ with $\rho:=\sqrt{\alpha^2+\beta^2}=\Delta\varphi\sqrt{D^2+s^2}$, and in terms of the original parameters
$$\widetilde L_{81}\simeq\Delta\varphi\Bigl\{4s+\frac{2D^2}{s}\Bigl[\ln\frac{4s}{D}-\frac56\Bigr]\Bigr\}.$$
Since $\widetilde L_{81}-4\Delta\varphi\,s$ is of order $\Delta\theta^2\ln\Delta\theta$, which dominates the $\Delta\theta^2$ term of the Taylor formula above, one obtains by adding the two the expansion of $L_{81}$ itself up to order $\Delta\theta^2$ on a logarithmic scale,
$$L_{81}=4\Delta\varphi\sin\theta_0+\frac{8\,\Delta\theta^2}{\Delta\varphi\sin\theta_0}\Bigl[\ln\frac{2\Delta\varphi\sin\theta_0}{\Delta\theta}-\frac56-\frac2{15}\Delta\varphi^2\sin^2\theta_0\Bigr]+o(\Delta\theta^2),$$
in which $\widetilde L_{81}$ already accounts for the logarithmic term and the constant $-5/6$, and the correction $E''(0)$ only shifts the constant.
For the values $\Delta\theta=6\pi/180$, $\Delta\varphi=50\pi/180$ and $\theta_0=65\pi/180$ one has $D=0.24$, $s\simeq0.90631$, $\alpha\simeq0.20944$ and $C\simeq0.065529<1/9$, so the lattice is rectangular, with $g_2\simeq0.54445$, $g_3\simeq-0.072409$, $e_r\simeq0.25542$, $\omega\simeq2.98311$, $\omega'\simeq1.97382\,i$, $\eta\simeq0.16719$, $a\simeq-1.24106\,i$ and $I(C)\simeq4.04218$. Hence $\widetilde L_{81}=4\alpha I(C)\simeq3.38637$, while numerical quadrature of the original integral gives $L_{81}\simeq3.37739$ (a relative difference of about $-2.7\cdot10^{-3}$), and the small-$D$ expansion gives $I\simeq4.0236$ and $\widetilde L_{81}\simeq3.3708$. 
The Taylor correction is $-\tfrac{16}{15}\Delta\varphi\,s\,\Delta\theta^2\simeq-0.00925$, so that $\widetilde L_{81}-\tfrac{16}{15}\Delta\varphi\,s\,\Delta\theta^2\simeq3.37711$ matches $L_{81}$ to $8\cdot10^{-5}$ in relative value.
%, whereas the log-scale expansion of $L_{81}$ above gives $3.3631$ (the remainder, which is only known to be $o(\Delta\theta^2)$, is not small at $C\simeq0.0655$).
